
\section{Data}\label{section:data}

The following analyses use data from multiple databases. This section provides information about the data used in this study and explains how the important variables were created.

\paragraph{Asset market data}

This study uses equity data from four sources: Global Financial Data (GFD), the Jorda-Schularik-Taylor Macrohistory Database (JST) mentioned in \citet*{Jorda2020}, IBES Global, and Factset. GFD offers two main stock return indices for each country. One shows the total return on stock exchanges in the country. The other shows the total return of all companies based in the country but listed on the London Stock Exchange.  

The primary variables of interest are both the dividend yield---defined as aggregate dividends over the calendar year divided by the price of the aggregate stock market index---and the annual excess return on the stock market. Excess returns are constructed assuming that investors have access to the same riskfree investment, in particular,  U.K. government bonds prior to 1914 and U.S. treasury bills after 1914.  This is because the returns on government bonds for the countries in my sample are generally not riskfree, and could be exposed to time-varying risks that equity assets are not exposed to \citep*{Miller2020}. Using home country government bonds may, therefore, erase part of the risk premium or induce measurement error in the dependent variable, reducing the statistical power of the results.

For all equity rate variables---like rates of return, dividend growth, and changes in dividend yields---I fill in missing observations in the GFD home stock market series using the JST data. Then I fill in missing observations using data from IBES Global, Factset, and the GFD data from the London Stock Exchange. Mixing these data sources gives an unbalanced panel data set of ex- and cum-dividend returns, dividend yields, and dividend growth over the longest time series possible for each country. For example, the data on dividend yield changes spans 201 years from 1816--2018 across 90 countries, with an average around \avgDivYldYears\ years of data for each country. However, because each series covers a different range, the number of observations varies throughout the paper. For more on how the asset pricing series are made, see Appendix~\ref{section:equity_app}.

\paragraph{Macroeconomic data}

Data on real GDP come from Maddison Historical Statistics, who use and expand upon data from \cite*{Barro2008} and provide the most comprehensive data available on these variables. Data on real consumption and the labor share of income come from the Penn World Tables.  These data are available from 1945 to the present. Data on income inequality come from the Standardized World Income Inequality Database (SWIID) who provide data on the Gini coefficient for up to 159 countries from 1960--2018. Finally, data on government revenue-GDP ratios come from GFD and data on tax revenue-GDP ratios come from the Relative Political Capacity Dataset. More information on the macroeconomic data used in the paper is provided in Appendix~\ref{section:macro_data}.

\paragraph{Political institutions data}

Data on political institutions come from the Varieties of Democracy (V-Dem) database.\footnote{This paper uses version 10.0 of the data.} V-Dem uses a team of over 3,500 country-specific experts to quantify levels of and trends in historical political institutions for most every country over the last two centuries. This allows them to provide the most detailed dataset possible to analyze changes in political institutions. V-Dem provides measures on both the level of electoral democratic institutions and other political outcomes. These other outcomes include the level and frequency of democratic protests, political violence, political polarization,  civil society activity, corruption, and bribery. More information on the measures used in the paper is provided in Appendix~\ref{section:vdem_ind}.

I also use measures on institutions from other sources where V-Dem does not provide data. These include the fraction of the population that is Catholic and a pro-competitive regulation index. These two measures come from the World Religion Project and the Fraser Institute's Economic Freedom Index. More detail on these series is provided in Appendix~\ref{section:otherpolinst}.

\paragraph{Events data}\label{section:eventsapp}
Data on events are primarily used as controls in the regressions below. Financial crises come from JST and \cite*{Reinhart2009} and are combined into a single financial crisis variable. Sovereign defaults also come from \cite*{Reinhart2009}. Recessions are taken from the GFD Dates database. Wars dates and locations come from the Correlates of War (CoW) data. International political crises come from the International Crisis Behavior (ICB) database as used in \cite*{Berkman2011}. Head of government and head of state deaths come from \cite*{Jones2009}, V-Dem, and Wikipedia. Data on head of government and head of state attempted assassinations also come from \cite*{Jones2009}. Regime changes are constructed using the regime information from V-Dem. More information on the events used in the paper can be found in Appendix~\ref{section:events_data}.

\subsection{Democratizations}\label{section:democratization_data}

Democratization and autocratization periods come from the Episodes of Regime Transformation (ERT) data.\footnote{This paper uses version 2.2 of the data.} There are two main advantages to using the ERT data. The first is that it is the only dataset to my knowledge that provides the start and end years of both democratization and autocratization episodes. Since asset prices are forward looking, this information is particularly important for this analysis. The second is that the ERT data provide detail on whether a democratization is sustained or reverts back to autocracy. For simplicity, I refer to these two potential outcomes as ``success'' or ``failure.'' By including both these types, I can avoid potential selection issues that come with conditioning on successful democratic transitions.

The ERT achieves this by examining changes in V-Dem's electoral democracy index above a certain threshold. This 0 to 1 index measures countries on the extent they embody the principles of electoral democracy. Countries that score highly generally respect principles of freedom of expression and association, have a high proportion of the population that can vote, and have elections that are competitive, clean, and fair. 

Since the asset pricing data are available prior to 1900, I extend the ERT data to back to 1816. To do this, I use the same procedure V-Dem uses to construct the post-1900 sample. This produces \demEpisodesPre\ additional democratization episodes for which asset pricing data are available. To obtain the latest possible end date for each democratization episode, I use data from \cite{Lindberg2018} to extend democratization episodes to their latest possible year.\footnote{\cite{Lindberg2018} follows a similar procedure to the ERT data, but with less conservative conditions on what constitutes the end of a democratization episode.} This gives \demYears\ democratization years across \demEpisodes\ episodes from 1816--2018 where I have dividend yield data.

Appendix~\ref{section:ERTdescription} provides more information on how the ERT data identifies democratizations and determines if they are successful or failed. Moreover, Appendix~\ref{section:case_studies} provides two case studies: one of the successful democratization in Sweden from 1917--1924 and the other of the failed democratization in France from 1847--1848. These case studies describe the historical background, asset pricing response, and subsequent redistribution (or lack thereof). An event timeline of all democratizations used for the asset pricing results is provided in Appendix Table~\ref{table:fastDem}.

\label{par:democratization_limit}
It is important to understand that these data---like all data on democratic institutions and democratizations---have their limits. Two particular caveats are worth noting. First, measures of democratization are noisy, which could introduce some degree of classical measurement error. If such measurement error were present, it would lead the results to be understated. Second, some have alleged that measures of democratic institutions come with some degree of left-wing ideological biases. This could lead them to rate more highly countries with greater amounts of redistribution, meaning left-leaning democratizations may be overrepresented in the analysis.  To address this, Appendix~\ref{section:robustness} shows the stock market results are quantitatively similar using 6 different measures of democratization.  This mitigates concerns that either noise or a particular set of ideological biases are driving the results.